The trenches Brienne had dug outside King’s Landing’s walls a week past were not the same trenches anymore, and she stood at their edge on the third morning of the city’s uncertain peace watching Garth’s men run the last clay pipe down into the eastern line, and thought that her father, who had taught her siege-craft from books because no one else had thought a girl worth teaching it properly, would not have recognized this particular use of it at all.

“Two more lengths of pipe and we’re finished the eastern run,” Garth told her, wiping soot from his young face with a sleeve already black with it. “Wildfire feeds down from the cisterns Her Grace had filled — forgive me, Ser Brienne, I keep calling her that, I don’t rightly know what to call a dead queen who isn’t queen of anything anymore — feeds down from the cisterns, along the pipe, into the trench itself. We light it at the head, here, where the ground’s driest, and the whole line goes up together rather than catching patchy the way an open trench fire generally does.”

“And the men standing in it when it goes.”

“There won’t be men in it, Ser Brienne. That’s the whole plan of the thing.” Garth said it with the particular carefulness of a young man repeating instructions he had memorized rather than fully trusted yet. “We fall back through the gaps Lord Tyrion’s marked, here and here, the moment the horns sound the retreat-line signal, and the dead follow into the trench because the dead don’t know a trap from a road, and then we light it from behind, once our own people are clear.”

“Once our own people are clear,” Brienne repeated, and found the phrase settling into her with more weight than Garth’s easy confidence seemed to carry for it. She had walked enough retreats this past year to know exactly how reliably a battle line held to its planned gaps once terror and blood and the particular chaos of ten thousand men fighting for their lives actually arrived to test the plan against reality.

\sceneBreak

Jaime found her there at midmorning, thinner than he’d been three days past and moving with the careful, deliberate economy of a man who had not properly slept since the morning he’d walked into his sister’s chamber, and Brienne watched him take in the pipe-work and the cistern and the whole elaborate machinery of the trap with the particular assessing look she remembered from Winterfell’s own defenses.

“It’s good work,” he said. “Better than I’d have managed myself, teaching Theon Greyjoy in a yard with half the time you’ve had.” He crouched at the trench’s edge, testing the depth with his golden hand against the packed clay wall. “Where’s the weak point.”

“There isn’t meant to be one.”

“There’s always one. I’ve built enough of these to know that much, whatever else this war’s taken from me to relearn it properly.” He looked up at her, and something in his face, tired as it was, held the old dry humor she’d come to trust more than his grimmer moods. “Show me where you’re worried, Brienne of Tarth. I’ve watched you long enough now to know you don’t stand at a trench’s edge with your arms crossed like that unless some part of you is already rehearsing how it goes wrong.”

She let out a breath, and pointed. “There. Where the trench bends round the old grain stores. We couldn’t run the pipe straight through — the stores are still half full, what’s left of the city’s grain, and Tyrion wouldn’t let us risk the wildfire that close to the one thing keeping ten thousand people fed for the next month. So the line kinks there, narrower than the rest of it, and narrower means slower falling back, if it comes to a real crush of men trying to clear it at once.”

“And slower falling back means someone’s still in the trench when it goes up.”

“That’s my worry, yes.” She said it plainly, because plain was the only way she knew how to say a true fear aloud, and found, saying it to Jaime rather than turning it over alone, that some of its weight eased simply in the telling. “I’ve moved the smallfolk labor gangs to the western stretch, away from that bend, for exactly that reason. But there’s still a company of Golden Company men holding that section under Ossric’s command, and Golden Company men are not mine to move without asking Roland Storm’s leave first, and Roland Storm has rather a great deal else occupying him just now.”

“Then ask him today. Not tomorrow.” Jaime rose, brushing clay from his knee, and looked out at the whole long snake of the trench-line stretching away toward the tree line where, in three days now, something considerably worse than an argument about troop placement was going to arrive to test every seam in this plan at once. “I’ve learned this past year that the thing which kills you in a battle like this one is never the enemy you planned for. It’s the gap nobody got round to closing because everyone assumed someone else had already seen to it.”

\sceneBreak

She found Roland Storm at the King’s Gate that same afternoon, overseeing a scorpion crew fitting dragonglass heads to bolts meant, until a week past, for nothing larger than an ordinary man, and watched him listen to her request with the particular weary patience of a captain who had spent the last three days being asked for something by nearly everyone in this castle.

“Ossric’s company can shift east,” he said, once she’d finished. “I’ll not pretend I’m glad to be moved a third time this week, Ser Brienne, but I understand a kinked trench better than most men standing on this wall, having watched enough of my own men die in worse ones in the Disputed Lands. Consider it done.”

“Thank you.”

“Don’t thank me yet. Thank me after, if there’s an after worth the thanking.” He turned back to the scorpion crew, watching a young Company armorer test the housing’s tension with more care than the girl seemed to think she had time to spare, and Brienne saw something pass across his face that she recognized from her own reflection some mornings this past week — the particular flat exhaustion of a person who had made too many decisions too quickly and would not know for certain, for days yet, which of them had been the right ones.

“Will they hold?” she asked. “The scorpions. Fitted for a dragon that isn’t the one they were built against.”

“Wisdom Hallyne’s apprentice swears the mechanism doesn’t care what’s on the other end of the bolt, so long as the housing holds the strain of a dragonglass head instead of steel, which is heavier and sits differently in the flight.” Roland’s mouth twisted, something between a grimace and dark amusement. “I find myself, Ser Brienne, in the peculiar position of trusting my life to a weapon my own employer’s dead queen built to kill a different dragon than the one it’s presently meant to kill, repurposed by a boy pyromancer nobody trusted with anything larger than a torch a month past. I don’t know if that’s grim comedy or simply what this whole war’s turned out to be, underneath all the songs anyone eventually writes about it.”

“Perhaps both,” Brienne said, and found, saying it, that she meant it more plainly than she’d expected to.

\sceneBreak

She walked the whole trench-line once more that evening, alone, the way she had the night before Winterfell’s own walls, torch in hand, measuring what could still be measured and making her uneasy peace with what could not. The grain-store kink had been widened as far as the stores themselves allowed, Ossric’s company shifted east as promised, the scorpion crews drilling their new dragonglass-fitted mechanisms until their arms ached from it. It was, she thought, as much as any army broken and starved and reduced to whatever this exhausted city could still muster was going to manage, in three days, against something that had already unmade a Wall seven hundred years standing and a castle older than any king currently claiming a throne.

\sceneBreak

Sansa came to walk the line herself that same afternoon, wrapped against the cold in furs that had, Brienne noted, begun to look considerably less like a lady’s traveling cloak and considerably more like the practical armor of a woman who had stopped dressing for anyone’s comfort but her own these past months.

“Lord Tyrion tells me the eastern kink is widened as far as it’ll go,” Sansa said, walking the trench’s edge beside her with the same careful, assessing attention she’d have once given a ledger. “And that Ossric’s men have shifted east of it.”

“They have. I’d not call the bend safe, my lady. I’d call it as safe as a bend in a wildfire trench is ever likely to be, which isn’t the same thing at all, whatever comfort either of us would rather take from it.”

“I didn’t come for comfort.” Sansa stopped, looking out over the whole long snake of pipe and packed earth stretching toward the dark tree line, and Brienne saw something settle in her face that was not quite peace, only the particular hard-won stillness of a woman who had learned to carry weight without it showing plainly on her. “I came because I wanted to see it with my own eyes before I have to stand behind it and trust it with people’s lives. I find I sleep better for having looked properly at the things I’m about to ask others to risk themselves on, rather than only reading them off Lord Tyrion’s chalk marks.”

“That’s more than most commanders trouble to do.”

“Most commanders didn’t grow up watching their mother trust the wrong men’s assurances,” Sansa said, quiet, and did not elaborate further, and Brienne did not ask her to. “It’ll hold, Ser Brienne. Or it won’t. But I’d rather know which, tonight, than be surprised by it three days from now with the whole city’s lives already spent on the not-knowing.”

She walked the rest of the line without further comment, asking a question here and there of the work-crews she passed, and Brienne watched her go with the particular quiet respect she had learned, across a hard year, to extend to people who did the unglamorous work of preparing properly rather than simply hoping loudly.

\sceneBreak

Grey Worm found her at the trench’s far end, near midnight, still and watchful the way he’d been every night since Winterfell, and Brienne understood, seeing him, that neither of them had come to this stretch of wall by accident so much as by the same restless inability to call the preparations finished while any daylight — or torchlight — remained to keep improving them.

“It will not be enough,” he said, not as despair, simply as arithmetic, the same flat accounting she’d heard from him digging trenches a week past. “I have counted what we have. I have counted, as near as any living person can, what is coming. The numbers do not balance kindly.”

“No,” Brienne agreed. “They don’t.” She looked out at the dark tree line, where nothing yet moved that torchlight could reach, though she understood, standing there, that somewhere beyond the edge of what any watch fire could show them, the gap that had been six days, then five, then four, was closing itself, patient and untiring, toward the last of whatever numbers either of them had left to count. “But we’ve built what we could build with what we had left to build it with. I don’t know that any of us gets more than that, in the end. I’ve stopped expecting the numbers to balance kindly, Grey Worm. I’ve only learned to make certain that whatever we do have counts for as much as it possibly can, when the counting finally starts.”

Grey Worm said nothing further, and the two of them stood a while longer at the trench’s edge in companionable silence, two soldiers keeping a private vigil neither one of them expected any reward for keeping, while behind them the city slept its last full night of not yet knowing, for certain, whether three days of frantic work was going to prove to have been enough.